\documentclass[11pt,a4paper]{article}

% ============================================================
% SCT Theory — NT-1b Phase 3: Combined SM Form Factors and c₁/c₂
% ============================================================

\usepackage[utf8]{inputenc}
\usepackage[T1]{fontenc}
\usepackage{amsmath,amssymb,amsfonts}
\usepackage[margin=2.5cm]{geometry}
\usepackage{hyperref}
\usepackage{booktabs}
\usepackage{xcolor}
\usepackage{fancyhdr}
\usepackage{enumitem}

\definecolor{sctblue}{RGB}{0,51,153}
\definecolor{verified}{RGB}{0,128,0}

\pagestyle{fancy}
\fancyhf{}
\fancyhead[L]{\small\textsc{SCT Theory --- Derivation}}
\fancyhead[R]{\small\textsc{NT-1b Phase 3}}
\fancyfoot[C]{\thepage}

\title{%
  \textbf{NT-1b Phase 3:}\\[6pt]
  \Large Combined SM Form Factors and $c_1/c_2$ Re-derivation\\[4pt]
  \large with $E = -R/4$ Correction and Full Particle Content
}
\author{David Alfyorov}
\date{March 2026}

\begin{document}
\maketitle


\begin{center}
\small\textit{Credits: Aliaksandr Samatyia contributed research-assistance and workflow support.}
\end{center}

\begin{center}
  \colorbox{verified!15}{\parbox{0.6\textwidth}{\centering
    \textbf{\color{verified} STATUS: VERIFIED}\\[2pt]
    \small independent cross-verification, multi-method checks, 354/354 checks PASS
  }}
\end{center}

\vspace{1em}

% ============================================================
\section{Overview}
% ============================================================

This document assembles the total Standard Model one-loop effective action
in the curvature-squared sector, using the verified form factors from
Phases 1 and 2 of NT-1b. We compute:

\begin{itemize}
  \item Total Weyl-squared form factor $\alpha_C(x)$ and its local limit $\alpha_C \equiv \alpha_C(0)$.
  \item Total $R^2$ form factor $\alpha_R(x,\xi)$ and its local limit $\alpha_R(\xi)$.
  \item The $c_1/c_2$ ratio in the $\{R^2, R_{\mu\nu}^2\}$ basis, including $\xi$-dependence.
  \item Scalar mode decoupling condition.
  \item UV asymptotic behavior of the total form factor.
\end{itemize}

% ============================================================
\section{Input: Verified Form Factors}
% ============================================================

From NT-1b Phases 1 and 2 (each verified to $\geq 98\%$ confidence via
multi-method verification), the individual spin-sector form factors are:

\subsection{Spin-0 (Real Scalar)}

\begin{align}
  h_C^{(0)}(x) &= \frac{1}{12x} + \frac{\varphi - 1}{2x^2}, &
  \beta_W^{(0)} &= h_C^{(0)}(0) = \frac{1}{120} \\
  h_R^{(0)}(x,\xi) &= f_{R,\mathrm{bis}}(x) + \xi\,f_{RU}(x) + \xi^2\,f_U(x), &
  \beta_R^{(0)}(\xi) &= h_R^{(0)}(0,\xi) = \frac{1}{2}\!\left(\xi - \frac{1}{6}\right)^{\!2}
\end{align}

\subsection{Spin-1/2 (4-Component Dirac Fermion)}

\begin{align}
  h_C^{(1/2)}(x) &= \frac{3\varphi - 1}{6x} + \frac{2(\varphi - 1)}{x^2}, &
  \beta_W^{(1/2)} &= -\frac{1}{20} \\
  h_R^{(1/2)}(x) &= \frac{3\varphi + 2}{36x} + \frac{5(\varphi - 1)}{6x^2}, &
  \beta_R^{(1/2)} &= 0
\end{align}

\subsection{Spin-1 (Massive Vector / Proca)}

\begin{align}
  h_C^{(1)}(x) &= \frac{\varphi}{4} + \frac{6\varphi - 5}{6x} + \frac{\varphi - 1}{x^2}, &
  \beta_W^{(1)} &= \frac{1}{10} \\
  h_R^{(1)}(x) &= -\frac{\varphi}{48} + \frac{11 - 6\varphi}{72x} + \frac{5(\varphi - 1)}{12x^2}, &
  \beta_R^{(1)} &= 0
\end{align}

Here $\varphi(x) = e^{-x/4}\sqrt{\pi/x}\,\mathrm{erfi}(\sqrt{x}/2)$
is the master function, with $\varphi(0) = 1$.

% ============================================================
\section{SM Particle Content and Counting Convention}
% ============================================================

\subsection{Standard Model Multiplicities}

\begin{center}
\renewcommand{\arraystretch}{1.3}
\begin{tabular}{lcc}
  \toprule
  \textbf{Species} & \textbf{Count} & \textbf{Notation} \\
  \midrule
  Real scalars (Higgs doublet $\to$ 4 real components) & 4 & $N_s = 4$ \\
  Two-component Weyl spinors & 45 & $N_f = 45$ \\
  Gauge bosons ($8 + 3 + 1$) & 12 & $N_v = 12$ \\
  \bottomrule
\end{tabular}
\end{center}

\subsection{Critical Convention: Weyl vs.\ Dirac Counting}

The form factors $h_C^{(1/2)}$ and $h_R^{(1/2)}$ are defined \textbf{per
4-component Dirac fermion}. The SM has $N_f = 45$ two-component Weyl spinors
(3 generations $\times$ 15 Weyl per generation). The number of Dirac-equivalent
fermions is:
\begin{equation}
  \boxed{N_D = \frac{N_f}{2} = \frac{45}{2} = 22.5}
  \label{eq:ND}
\end{equation}

This follows from the fact that each 4-component Dirac spinor contains 2 Weyl
spinors. The half-integer count is algebraically consistent.

\textbf{Reference:} Codello, Percacci, Rahmede (CPR), 0805.2909, eq.~(3.14),
which explicitly uses $n_D$ for the number of 4-component Dirac fermions.

\textbf{Bug fixed:} Prior to Phase 3, the code used $N_f \cdot h_C^{(1/2)}$
(treating $N_f=45$ as the Dirac count), which overestimated the fermionic
contribution by a factor of 2. This was corrected to $(N_f/2) \cdot h_C^{(1/2)}$.

% ============================================================
\section{Total Weyl Form Factor $\alpha_C$}
% ============================================================

\subsection{Full Nonlocal Form Factor}

The total Weyl-squared form factor is:
\begin{equation}
  \alpha_C(x) = N_s\,h_C^{(0)}(x) + N_D\,h_C^{(1/2)}(x) + N_v\,h_C^{(1)}(x)
  \label{eq:alphaC-full}
\end{equation}

\subsection{Local Limit}

Setting $x = 0$ (using the verified $\beta$ values):
\begin{align}
  \alpha_C &= N_s \cdot \beta_W^{(0)} + N_D \cdot \beta_W^{(1/2)} + N_v \cdot \beta_W^{(1)} \nonumber \\
  &= 4 \cdot \frac{1}{120} + \frac{45}{2} \cdot \left(-\frac{1}{20}\right) + 12 \cdot \frac{1}{10} \nonumber \\
  &= \frac{4}{120} - \frac{45}{40} + \frac{12}{10} \nonumber \\
  &= \frac{4}{120} - \frac{135}{120} + \frac{144}{120} \nonumber \\
  &= \frac{4 - 135 + 144}{120}
\end{align}

\begin{equation}
  \boxed{\alpha_C = \frac{13}{120}}
  \label{eq:alphaC}
\end{equation}

This is \textbf{positive} and independent of $\xi$. The positivity means the
massive spin-2 mode in linearized higher-derivative gravity is a ghost
(negative residue) --- a well-known feature of quadratic gravity (Stelle, 1977).

\subsection{Sign Structure}

\begin{center}
\renewcommand{\arraystretch}{1.2}
\begin{tabular}{lccc}
  \toprule
  \textbf{Sector} & \textbf{$N_i \cdot \beta_W^{(s_i)}$} & \textbf{Value} & \textbf{Sign} \\
  \midrule
  Scalar ($s=0$)     & $4/120$       & $1/30$   & $+$ \\
  Fermion ($s=1/2$)  & $-45/40$      & $-9/8$   & $-$ \\
  Vector ($s=1$)     & $12/10$       & $6/5$    & $+$ \\
  \midrule
  \textbf{Total}     & $13/120$      & $\approx 0.1083$ & $+$ \\
  \bottomrule
\end{tabular}
\end{center}

The fermionic sector opposes the scalar and vector sectors, but the combined
scalar+vector contribution ($4/120 + 144/120 = 148/120$) outweighs the
fermionic contribution ($-135/120$) by $13/120$.

% ============================================================
\section{Total $R^2$ Form Factor $\alpha_R(\xi)$}
% ============================================================

\subsection{Local Limit}

\begin{align}
  \alpha_R(\xi) &= N_s \cdot \beta_R^{(0)}(\xi) + N_D \cdot \beta_R^{(1/2)} + N_v \cdot \beta_R^{(1)} \nonumber \\
  &= 4 \cdot \frac{1}{2}\!\left(\xi - \frac{1}{6}\right)^{\!2} + \frac{45}{2} \cdot 0 + 12 \cdot 0
\end{align}

\begin{equation}
  \boxed{\alpha_R(\xi) = 2\!\left(\xi - \frac{1}{6}\right)^{\!2}}
  \label{eq:alphaR}
\end{equation}

\begin{itemize}
  \item $\alpha_R \geq 0$ for all $\xi$ (perfect square).
  \item $\alpha_R(1/6) = 0$ (conformal coupling).
  \item $\alpha_R(0) = 2 \cdot (1/6)^2 = 1/18$ (minimal coupling).
  \item Only the scalar sector contributes; Dirac and vector $\beta_R$ vanish
        by conformal invariance of spinor and Maxwell fields.
\end{itemize}

% ============================================================
\section{Basis Conversion and $c_1/c_2$ Ratio}
% ============================================================

\subsection{From $\{C^2, R^2\}$ to $\{R^2, R_{\mu\nu}^2\}$}

The Weyl tensor squared in 4D, using the Gauss--Bonnet identity
($R_{\mu\nu\rho\sigma}^2 = 4R_{\mu\nu}^2 - R^2 + E_4$, with $E_4$ topological):
\begin{equation}
  C_{\mu\nu\rho\sigma}C^{\mu\nu\rho\sigma} = 2\,R_{\mu\nu}R^{\mu\nu} - \frac{2}{3}\,R^2
  \quad \text{(mod $E_4$)}
\end{equation}

Therefore:
\begin{equation}
  \alpha_C\,C^2 + \alpha_R\,R^2
  = 2\alpha_C\,R_{\mu\nu}^2 + \left(\alpha_R - \frac{2}{3}\alpha_C\right)R^2
\end{equation}

With our convention $S = \int\!\sqrt{g}\,(c_1\,R^2 + c_2\,R_{\mu\nu}^2)$:
\begin{equation}
  c_1 = \alpha_R - \frac{2}{3}\alpha_C, \qquad c_2 = 2\alpha_C
  \label{eq:c1c2-def}
\end{equation}

\subsection{The Ratio}

\begin{align}
  \frac{c_1}{c_2} &= \frac{\alpha_R - \frac{2}{3}\alpha_C}{2\alpha_C}
  = -\frac{1}{3} + \frac{\alpha_R}{2\alpha_C}
\end{align}

Substituting $\alpha_C = 13/120$ and $\alpha_R(\xi) = 2(\xi - 1/6)^2$:

\begin{equation}
  \boxed{
    \frac{c_1}{c_2} = -\frac{1}{3} + \frac{120\!\left(\xi - \tfrac{1}{6}\right)^{\!2}}{13}
  }
  \label{eq:c1c2-ratio}
\end{equation}

\subsection{Special Cases}

\begin{center}
\renewcommand{\arraystretch}{1.3}
\begin{tabular}{lccl}
  \toprule
  \textbf{Coupling} & $\xi$ & $c_1/c_2$ & \textbf{Note} \\
  \midrule
  Conformal & $1/6$ & $-1/3$ & Reproduces pure-Dirac prediction \\
  Minimal   & $0$   & $-1/13$ & Small negative ratio \\
  $\xi = 1$ &       & $-1/3 + 120 \cdot (5/6)^2/13 = 79/13 \approx 6.08$ & Positive ratio (unusual) \\
  \bottomrule
\end{tabular}
\end{center}

At $\xi = 1/6$, we recover the result $c_1/c_2 = -1/3$ originally derived
from the pure Dirac operator (prediction\_01), where $\alpha_R = 0$ and only
$\alpha_C$ contributes. The full SM derivation shows this is a special case
of the more general $\xi$-dependent formula.

% ============================================================
\section{Scalar Mode Decoupling}
% ============================================================

The massive spin-0 gravitational mode in quadratic gravity has mass
$m_0^2 \propto 1/(3c_1 + c_2)$. Computing:

\begin{align}
  3c_1 + c_2 &= 3\!\left(\alpha_R - \frac{2}{3}\alpha_C\right) + 2\alpha_C \nonumber \\
  &= 3\alpha_R - 2\alpha_C + 2\alpha_C \nonumber \\
  &= 3\alpha_R
\end{align}

\begin{equation}
  \boxed{3c_1 + c_2 = 3\alpha_R(\xi) = 6\!\left(\xi - \frac{1}{6}\right)^{\!2}}
  \label{eq:scalar-mode}
\end{equation}

The scalar mode decouples ($m_0^2 \to \infty$) if and only if
$3c_1 + c_2 = 0$, which requires:

\begin{equation}
  \xi = \frac{1}{6} \qquad \text{(conformal coupling of the Higgs field)}
\end{equation}

\textbf{Physical interpretation:} The cancellation of $\alpha_C$ in $3c_1 + c_2$
is a structural consequence of the Gauss--Bonnet identity. The scalar graviton
sector is controlled entirely by $\alpha_R$, which vanishes at conformal coupling.
This provides a dynamical mechanism for scalar-mode avoidance: if the Higgs
field is conformally coupled, the problematic scalar graviton is absent from
the spectrum.

% ============================================================
\section{UV Asymptotic Behavior}
% ============================================================

For $x \to \infty$, $\varphi(x) \to 2/x$ (from the asymptotic expansion of
$\mathrm{erfi}$). The leading $1/x$ behavior of each form factor:

\begin{align}
  x \cdot h_C^{(0)}(x) &\to \frac{1}{12} \\
  x \cdot h_C^{(1/2)}(x) &\to -\frac{1}{6} \\
  x \cdot h_C^{(1)}(x) &\to -\frac{1}{3}
\end{align}

Note: for the vector, the $\varphi/4 \to 2/(4x) = 1/(2x)$ term contributes
at leading order, combining with $(6\varphi-5)/(6x) \to -5/(6x)$ to give
$1/(2x) - 5/(6x) = -1/(3x)$. This was verified numerically and caught an
error in the D-R stage derivation.

The total UV coefficient:
\begin{align}
  \lim_{x\to\infty} x \cdot \alpha_C(x) &= N_s \cdot \frac{1}{12} + N_D \cdot \left(-\frac{1}{6}\right) + N_v \cdot \left(-\frac{1}{3}\right) \nonumber \\
  &= \frac{4}{12} - \frac{45}{12} - \frac{48}{12}
\end{align}

\begin{equation}
  \boxed{\lim_{x\to\infty} x \cdot \alpha_C(x) = -\frac{89}{12}}
  \label{eq:uv-asymptotic}
\end{equation}

The total form factor $\alpha_C(x)$ starts positive at $x=0$ ($13/120 > 0$)
but becomes negative for large $x$ (the $-89/12$ coefficient is negative).
Therefore $\alpha_C(x)$ crosses zero at some finite $x_0$, indicating a sign
change of the effective Weyl-squared coupling in the UV.

% ============================================================
\section{Complete Phase 3 Effective Action}
% ============================================================

The one-loop curvature-squared effective action for the full SM is:
\begin{equation}
  \Gamma^{(1)}_{\mathrm{curv}} = \frac{1}{16\pi^2}\int d^4x\,\sqrt{g}\left[
    \alpha_C\!\left(\frac{\Box}{\Lambda^2}\right)\!C_{\mu\nu\rho\sigma}C^{\mu\nu\rho\sigma}
    + \alpha_R\!\left(\frac{\Box}{\Lambda^2},\xi\right)\!R^2
  \right]
\end{equation}

with the local limits:
\begin{align}
  \alpha_C &= \frac{13}{120}, &
  \alpha_R(\xi) &= 2\!\left(\xi - \frac{1}{6}\right)^{\!2}
\end{align}

In the $\{R^2, R_{\mu\nu}^2\}$ basis:
\begin{align}
  c_1 &= \alpha_R - \frac{2}{3}\alpha_C = 2\!\left(\xi - \frac{1}{6}\right)^{\!2} - \frac{13}{180} \\
  c_2 &= 2\alpha_C = \frac{13}{60}
\end{align}

% ============================================================
\section{Summary of Verified Results}
% ============================================================

\begin{center}
\renewcommand{\arraystretch}{1.4}
\begin{tabular}{lcc}
  \toprule
  \textbf{Quantity} & \textbf{Value} & \textbf{Status} \\
  \midrule
  $\alpha_C$ & $13/120$ & \textcolor{verified}{VERIFIED} \\
  $\alpha_R(\xi)$ & $2(\xi - 1/6)^2$ & \textcolor{verified}{VERIFIED} \\
  $c_1/c_2$ (general) & $-1/3 + 120(\xi-1/6)^2/13$ & \textcolor{verified}{VERIFIED} \\
  $c_1/c_2|_{\xi=1/6}$ & $-1/3$ & \textcolor{verified}{VERIFIED} \\
  $c_1/c_2|_{\xi=0}$ & $-1/13$ & \textcolor{verified}{VERIFIED} \\
  $3c_1 + c_2$ & $6(\xi - 1/6)^2$ & \textcolor{verified}{VERIFIED} \\
  UV asymptotic ($x \cdot \alpha_C$) & $-89/12$ & \textcolor{verified}{VERIFIED} \\
  Scalar decoupling condition & $\xi = 1/6$ & \textcolor{verified}{VERIFIED} \\
  \bottomrule
\end{tabular}
\end{center}

% ============================================================
\section{Verification Summary}
% ============================================================

\begin{center}
\renewcommand{\arraystretch}{1.2}
\begin{tabular}{lcc}
  \toprule
  \textbf{Layer} & \textbf{Result} & \textbf{Checks} \\
  \midrule
  1. Analytic (limits, signs, dimensions) & PASS & 13/13 \\
  2. Numerical (mpmath, 120 digits) & PASS & 8/8 \\
  2.5. Randomized property testing (hypothesis) & PASS & 304/304 \\
  3. Literature (CPR, Avramidi, Stelle) & PASS & 8/8 \\
  4. Independent dual derivation & PASS & 5/5 \\
  4.5. Independent CAS verification (Fraction, mpmath, SymPy) & PASS & 5/5 \\
  5--6. Formal (rational arithmetic) & PASS & 11/11 \\
  \midrule
  \textbf{Total} & \textbf{PASS} & \textbf{354/354} \\
  \bottomrule
\end{tabular}
\end{center}

One error caught during verification: D-R stage computed the UV asymptotic
as $-161/12$ (using $\varphi \to 0$ instead of $\varphi \to 2/x$). Error
confined to UV asymptotic only; all other results independently confirmed.

Full test suite: 2046/2046 passed, 0 failed.

\vspace{1em}
\noindent
\textit{This document is part of the SCT Theory derivation archive.
Phase 3 completes the NT-1b program for the full SM curvature-squared sector.}

\end{document}
